\[g(x) = 2x - x^{2}.\]

\smallskip

\begin{enumerate}
\item 
Le trinôme $g(x)$ est dérivable sur $\R$, donc sur $[0~;~1]$ et sur cet intervalle :

$g'(x) = 2 - 2x = 2(1 - x)$.

Or on sait que $0\leqslant x \leqslant 1\Rightarrow - x \leqslant 0 \leqslant 1 - x$ ou en lisant de droite à gauche 

$1 - x \geqslant 0 \Rightarrow 2(1 - x) \geqslant 0 \iff g'(x) \geqslant 0$ : la dérivée étant positive sur $[0~;~1]$ et ne s'annulant qu'en $x = 1$, la fonction $g$ est strictement croissante de $g(0 ) = 0$ à $g(1) = 2 - 1^2 = 1$.
\end{enumerate}

On considère la suite $\left(u_{n}\right)$ définie par $\left\{\begin{array}{l c l}u_{0}&=&\dfrac{1}{2}\\ u_{n+1}&=&g\left(u_{n}\right)\end{array}\right.$ pour tout entier naturel $n$.

\begin{enumerate}
\setcounter{enumi}{1}
\item 
$\bullet~~u_1 = 2 \times \dfrac{1}{2} - \left(\dfrac{1}{2}\right)^2 = 1 - \dfrac14 = \dfrac34 = 0,75$ ;

$\bullet~~u_2 = 2 \times \dfrac34 - \left(\dfrac34\right)^2 = \dfrac32 - \dfrac{9}{16} = \dfrac{24}{16} -  \dfrac{9}{16} = \dfrac{15}{16} = \np{0,9375}$.
\item 
Démonstration par récurrence :

$\bullet~~$\emph{Initialisation} :

On a $0 < \dfrac12 < \dfrac34 < 1$, soit $0 < u_0 < u_1 < 1$ : l'encadrement est vrai au rang zéro.

$\bullet~~$\emph{Hérédité} :
Supposons que pour $n \in \N$, on ait $0 < u_{n} < u_{n+1}< 1$.

Alors par stricte croissance sur [0~;~1] de la fonction $g$, on a :

$g(0) < g\left(u_n\right) < g\left(u_{n+1} \right) < g(1)$, soit d'après les résultats précédents :

$0 < u_{n+1} < u_{n+2} < 1$ : l'encadrement est encore vrai au rang $n + 1$.

\emph{Conclusion} : la relation est vraie au rang $n = 0$ et si elle est vraie au rang $n \in \N$ elle l'est encore au rang $n + 1$ : par le principe de récurrence, on a donc 

Pour tout entier naturel $n$,\:: $0 < u_{n} < u_{n+1} < 1$.
\item 
Le résultat précédent montre que la suite $\left(u_{n}\right)$ est croissante  et qu'elle est majorée par 1 : elle converge donc vers une limite $\ell \leqslant 1$.
\item 
La relation $g\left(u_n\right) = u_{n+1} = 2u_n - u_n^2$ donne à la limite car $g$ est continue car dérivable sur l'intervalle [0~;~1] : $\ell = 2\ell - \ell^2 \iff $

$\ell - \ell^2 = 0 \iff \ell (1 - \ell) = 0$ soit 

$\left\{\begin{array}{l c l}
\ell&=&0 \:\:\text{ou}\\
1 - \ell&=&0
\end{array}\right. \iff \left\{\begin{array}{l c l}
\ell&=&0 \:\:\text{ou}\\
1  &=&\ell
\end{array}\right.$

La solution $\ell = 0 $ n'est pas possible (la suite est croissante ) ; il reste $\displaystyle\lim_{n \to + \infty} u_n = 1$
\end{enumerate}

On considère la suite $\left(v_{n}\right)$ définie pour tout entier naturel $n$ par $v_{n}=\ln \left(1 - u_{n}\right)$.

\begin{enumerate}
\setcounter{enumi}{5}
\item 
On a pour $n \in \N, \: v_{n+1} = \ln \left(1 - u_{n+1}\right) = \ln \left[1 - \left(2u_n - u^2_n \right) \right] = \ln \left(1 - 2u_n + u_n^2 \right) =$

$ \ln \left( 1 - u_n\right)^2 = 2\ln \left( 1 - u_n\right)$ (car $1 - u_n > 0$ voir la récurrence ci-dessus, donc $\ln \left( 1 - u_n\right)$ existe). Or $\ln \left(1 - u_n\right) = v_n$.

Finalement : $v_{n + 1} = 2v_n$ ce qui montre que la suite  $\left(v_{n}\right)$ est une suite géométrique de raison 2 et de premier terme $v_0 = \ln \left(1 - u_0\right) = \ln \left(1 - \dfrac12\right) = \ln \dfrac12 = - \ln 2$.
\item 
On sait qu'alors pour tout $n \in \N, \:v_n = v_0 \times 2^n$, soit $v_n = - \ln 2 \times 2^n$.
\item 
La relation $v_{n}=\ln \left(1 - u_{n}\right)$ donne donc :

$- \ln 2 \times 2^n = \ln \left(1 - u_{n}\right) \iff \e^{- \ln 2 \times 2^n} = \e^{\ln \left(1 - u_{n}\right)}$, (par croissance de la fonction exponentielle), soit encore :

$\e^{- \ln 2 \times 2^n} = 1 - u_n \iff u_n = 1 - \e^{- \ln 2 \times 2^n}$.

Or on sait que $\displaystyle\lim_{n \to + \infty}- \ln 2 \times 2^n = - \infty$, donc $\displaystyle\lim_{n \to + \infty} \e^{- \ln 2 \times 2^n} = 0$ et par conséquent :
\[\displaystyle\lim_{n \to + \infty} u_n = 1.\]

\item 
\end{enumerate}

\begin{center}
\begin{tabular}{|>{\texttt} l|}\hline

\textbf{def }seuil():\\
\qquad n=0\\
\qquad u=0.5\\
\qquad \textbf{while} u < 0.95 :\\
\qquad \qquad n={\red n + 1}\\
\qquad \qquad {\red u=2*u - u**2}\\
\qquad \textbf{return} n\\ \hline
\end{tabular}

\end{center}
